\documentclass[10pt,a4paper]{amsart}
\usepackage[margin=19mm]{geometry}
\usepackage{amsmath,amssymb,mathtools}
\usepackage[scr=rsfs,cal=euler]{mathalfa}
\usepackage{xfrac}
\usepackage[unicode,hidelinks,bookmarks=false]{hyperref}
\newcommand{\ie}{i.e.\ }
\newcommand{\cf}{cf.\ }
\newcommand{\resp}{respectively\ }
\newcommand{\Subsection}{\subsection}
\providecommand{\nobreakdash}{\nobreak\hbox{-}}
\setlength{\emergencystretch}{2em}
\multlinegap0pt
\usepackage{bm}
\usepackage{bbm}
\usepackage{dsfont}
\usepackage{xspace}
\usepackage{cancel}
\usepackage{mathtools}
\usepackage{amsthm}
\makeatletter
\@ifundefined{Proposition}{\newtheorem{Proposition}{Proposition}}{}
\makeatother
\makeatletter
\expandafter\ifx\csname Figure\endcsname\relax
\newenvironment{Figure}
  {\par\medskip\noindent\minipage{\linewidth}}
  {\endminipage\par\medskip}
\fi
\makeatother
\title[On a Variation of Gambler's Ruin Problem]{On a Variation of Gambler's Ruin Problem}
\author{Zhiyi Chi, Vladimir Pozdnyakov}
\date{Source: arXiv:2506.00990v1 (CC BY 4.0)}
\begin{document}
\maketitle

\noindent\emph{Source and licence.} On a Variation of Gambler's Ruin Problem. Version arXiv:2506.00990v1 (CC BY 4.0), distributed under the Creative Commons Attribution 4.0 International licence, as recorded in the arXiv licence metadata for this exact version and shown on the abstract page https://arxiv.org/abs/2506.00990v1. Full rights notice: licenses/arxiv-2506.00990-CC-BY-4.0.txt.\par\smallskip

\noindent\textit{Editorial scope.} Counted unit: the Proposition whose source label is \texttt{None} in the article (source0.tex lines 539--546, source label \texttt{None}), delivered together with the article's own complete original proof of it (source0.tex lines 547--573, one \texttt{proof} environment of 27 lines). No mathematics is written, repaired or re-derived: statement and proof are the article's own text line for line. Required context retained uncounted: none. Every definition, display and label of those lines is retained unchanged. Omitted from this package and not redistributed: 1--538, 574--606. Nothing inside a retained excerpt is omitted or altered: every retained body is delivered exactly as the article's lines are written, so no omission receipt of any kind accompanies this package. Citations: the retained excerpts carry no citation command at all, so no bibliography entry of the article is transcribed and no reference is redistributed. Reference adaptation: none was needed and none was made; every reference inside the delivered unit resolves inside it, so no unresolved reference or citation is produced. Printed slips kept literal: no printed word, symbol, hypothesis or number of the counted statement or of its proof was altered. Layout and class: the article's own class is \texttt{amsart} with packages including natbib, color, graphicx, float, multicol, amsmath, fullpage, booktabs, amsthm, hyperref, amsfonts, bm, geometry, setspace; the wrapper is the lane's pinned \texttt{amsart} editorial wrapper at 10pt on A4, which restates the theorem-like environments the retained text needs and adds the article's own macro definitions that the retained text needs (none), with extra wrapper packages (bm, bbm, dsfont, xspace, cancel, mathtools). The delivered numbering is the unit's own. Assets: no retained span contains a figure environment, a graphics-inclusion command, a PDF-page inclusion, a table or a TikZ picture; no image file is copied into this package, the package declares no asset dependency and no image file is redistributed with it. Licence: version arXiv:2506.00990v1 of this article is distributed under the Creative Commons Attribution 4.0 International licence, as recorded in the arXiv licence metadata for this exact version and shown on the abstract page https://arxiv.org/abs/2506.00990v1. A full rights notice is delivered at licenses/arxiv-2506.00990-CC-BY-4.0.txt. Characters: every retained excerpt is pure ASCII, so the delivered engine, XeLaTeX with its default Latin Modern text font, reports no missing character.
\begin{Proposition}
Let the Markov chain $\{Y_k\}_{k\geq 1}$ with state set $\Delta=\{0,1\}$ be irreducible and
aperiodic (or, equivalently, its transition matrix is strictly positive).
Then the Markov chain $\{X_k\}_{k\geq 1}$ is periodic if and only if
we have one of the following combinations of words: $U = y^t(1-y)^s$ and
$D = (1-y)^sy^t$, where $y \in \Delta$, $t, s \geq 1$, but either $s = 1$ or
$t = 1$. Here, $y^t$ denotes a word formed by repeating the letter $y$ $t$ times.
\end{Proposition}

\begin{proof}
Assume that $U$ is always followed by $D$, and $D$ is always followed by $U$.
First, it is obvious that neither $U$ nor $D$ are runs; both must include both
the letters $0$ and $1$. Assume the last letter in $D$ is $1$. Consider the following
possible text: $U0^nU$, where $n$ is greater than the maximum of  lengths of $D$ and $U$.
Since $U$ can only occur at the beginning and the end of the text $U0^nU$,
$D$ must occur in the middle. Moreover, the last letter of $D$, which is $1$,
must be in the second $U$ at the end.  This means that $D$ starts with a run of $0$
(with a length of at least one).

Next, let us consider another possible text: $U0^n1^nU$. If $U=0^i1^j$ (that is, there
is a $U$ in the middle of $U0^n1^nU$), then by considering $U0^nU$, we get that $D=0^t1^s$,
where $t, s \geq 1$. If $U\neq 0^i1^j$, then there are only two $U$s in $U0^n1^nU$, and the
last $1$ of $D$ must be in $1^nU$. Therefore, again $D=0^t1^s$, where $t, s \geq 1$.
Using the same argument, we get that $U=0^i1^j$ or $U=1^i0^j$, where $i, j \geq 1$.
But by considering the text $DD$, we see that $U=0^i1^j$ must be excluded because
it cannot occur in $DD$ (remember that $U$ is not a subword of $D$). Thus, $U=1^i0^j$. If $t<i$,
then $U$ cannot occur in $DD$. Therefore, $t \geq i$.
By considering the occurrence of $D$ in $UU$, we get that $t=i$.
Using similar arguments, we also have $s=j$, that is, $D=0^s1^t$ and $U=1^t0^s$.

Finally, if both $s,t>1$, consider the text $D(01)^nD$ and observe that $U$ cannot
occur in this text. Therefore, either $s=1$ or $t=1$ (or both). A direct check
shows that if $D=0^s1^t$ and $U=1^t0^s$, with $t,s\geq 1$, either $s = 1$ or
$t = 1$, then $U$ and $D$ alternates in any text generated by $\{Y_k\}_{k\geq 1}$.
The other two cases are obtained by interchanging $0$ and $1$.
\end{proof}

\end{document}
